\ifx\allfiles\undefined
\documentclass[12pt, a4paper, oneside, UTF8]{ctexbook}
\def\path{../config}
\input{\path/package.tex}

% 定理环境设置
\input{\path/theorem1_zh.tex}

\input{\path/custom.tex}

% 修改参数改变封面样式，0 默认原始封面、内置其他1、2、3种封面样式
\def\myIndex{3}


\ifnum\myIndex>0
    \input{\path/cover_package_\myIndex} 
\fi

\def\myTitle{高等数学笔记}
\def\myAuthor{M.K.}
\def\myDateCover{\today}
\def\myDateForeword{\today}
\def\myForeword{前言}
\def\myForewordText{
    
    这是一个基于\LaTeX{}模板的高数笔记．
    内容参考了包括但不限于以下资料：
    \begin{itemize}
        \item 《高等数学》（同济大学出版社）第七版
        \item 《吉米多维奇高等数学习题精选精讲》（山东科学技术出版社）第二版
        \item 《高等数学作业集》（哈尔滨工业大学出版社）第一版
        \item \href{https://space.bilibili.com/1035929235}{B站UP主：一高数}的高数课程
        \item \href{https://space.bilibili.com/42428180}{B站UP主：考研竞赛数学}的试题讲解与高数课程
        \item \href{https://space.bilibili.com/391930545}{B站UP主：Maki的完美算数教室}的高数内容

    \end{itemize}

    封面来源于\href{https://www.pixiv.net/users/3952}{おにねこ}的作品\href{https://www.pixiv.net/artworks/136551599}{雨雫とプレアデス}
    如有侵权请联系我删除．
}
\def\mySubheading{副标题}


\begin{document}
% \input{../config/cover}
\else
\fi
\chapter{泰勒公式}
\section{$f(x)$的展开式}
\begin{example}
    求函数$f(x)=x^2 \ln x $在点$x=1$处带有拉格朗日余项的二阶泰勒展开式
\end{example}
\begin{solution}
    由$f(x)=x^2 \ln x$，可得
    \[
        f'(x)=2x\ln x + x, \quad f''(x)=2\ln x + 3, \quad f'''(x)=\frac{2}{x}
    \]
    因此在$x=1$处，有
    \[
        f(1)=0, \quad f'(1)=1, \quad f''(1)=3, \quad f'''(1)=2
    \]
    由泰勒公式知，$f(x)$在$x=1$处带有拉格朗日余项的二阶泰勒展开式为
    \[
        f(x)=0 + 1(x-1) + \frac{3}{2!}(x-1)^2 + \frac{2}{3!}\xi (x-1)^3,
    \]
    其中$\xi$介于$x$与$1$之间．
\end{solution}
\begin{example}
    求函数$f(x)=xe^x$带有佩亚诺余项的n阶麦克劳林展开式
\end{example}
\begin{solution}
    法一：
    由$e^x$的麦克劳林展开式知
    \[
        e^x=1 + \frac{x^1}{1!} + \frac{x^2}{2!} + \cdots + \frac{x^n}{n!} + o(x^n)
    \]
    因此
    \[
        f(x)=xe^x=x\left(1 + \frac{x^1}{1!} + \frac{x^2}{2!} + \cdots + \frac{x^n}{n!} + o(x^n)\right)
    \]
    即
    \[
        f(x)=x + x^2 + \dfrac{x^3}{2!} + \cdots + \frac{x^n}{(n-1)!} + o(x^n)
    \]
    法二：
    由$f(x)=xe^x$，可得
    \begin{align*}
        f^{(n)}(x)&=(xe^x)^{(n)}=x(e^x)^{(n)} + \binom{n}{1}(x)'(e^x)^{(n-1)} + \cdots + \binom{n}{n-1}(x)^{(n-1)}(e^x)' + \binom{n}{n}(x)^{(n)}(e^x)\\
        &=xe^x + n e^x + 0 + \cdots + 0 = (x+n)e^x
    \end{align*}
    因此在$x=0$处，有
    \[
        f^{(n)}(0)=n
    \]
    由麦克劳林公式知，$f(x)$带有佩亚诺余项的n阶麦克劳林展开式为
    \[
        f(x)=x + x^2 + \dfrac{x^3}{2!} + \cdots + \frac{x^n}{(n-1)!} + o(x^n)
    \]
\end{solution}
\section{泰勒公式求极限}
\begin{example}
    求极限$\lim\limits_{x \to 0} \dfrac{\sin x-x\cos x}{x^3}$
\end{example}
\begin{solution}
    \begin{align*}
        \lim_{x \to 0} \dfrac{\sin x-x\cos x}{x^3}&=\lim_{x \to 0} \dfrac{\left(x - \dfrac{x^3}{3!} + o(x^3)\right) - x\left(1 - \dfrac{x^2}{2!} + o(x^2)\right)}{x^3}\\
        &=\lim_{x \to 0} \dfrac{\dfrac{1}{3}x^3 + o(x^3)}{x^3}
        = \dfrac{1}{3}
    \end{align*}
    \myRemark{
    \textcolor{red}{$x$最高次数即为泰勒展开的阶数．}
}
\end{solution}
\section{泰勒公式求高阶导}
\begin{example}
    求函数$f(x)=x^2 \ln(1+x)$在$x=0$处的$n$阶导数$f^{(n)}(0),(n \geq 3)$
\end{example}
\begin{solution}
    由$\ln(1+x)$的麦克劳林展开式知
    \[
        \ln(1+x)=x - \frac{x^2}{2} + \frac{x^3}{3} - \frac{x^4}{4} + \cdots + (-1)^{n-1}\frac{x^n}{n} + R_n (x)
    \]
    因此
    \begin{align*}
        f(x)&=x^2 \ln(1+x)\\
        &=x^2\left[x - \frac{x^2}{2} + \frac{x^3}{3} - \frac{x^4}{4} + \cdots + (-1)^{n-1}\frac{x^n}{n} + R_n (x)\right]\\
        &=x^3 - \frac{x^4}{2} + \frac{x^5}{3} - \frac{x^6}{4} + \cdots + (-1)^{n-3}\frac{x^{n}}{n-2} + R_n (x)
    \end{align*}
    由泰勒展开的唯一性，可得
    \[
        \dfrac{f^{(n)}(0)}{n!}=\dfrac{(-1)^{n-3}}{n-2}
    \]
    即
    \[
        f^{(n)}(0)=(-1)^{n-3}\dfrac{n!}{n-2}, \quad (n \geq 3)
    \]
\end{solution}
\section{利用泰勒公式证明}
\begin{example}
    函数$f(x)$在$[0,1]$上二阶可导，且$f(0)=1,f'(0)=0,f''(x) \le 2$.\\
    证明：$\max\limits_{x \in [0,1]} f(x) \le 2$
\end{example}
\begin{proof}
    由泰勒公式知，存在$\xi$介于$x$与$0$之间，使得
    \[
        f(x)=f(0) + f'(0)x + \frac{f''(\xi)}{2}x^2
    \]
    因为$f(0)=1,f'(0)=0,f''(x) \le 2$，所以
    \[
        f(x) \le 1 + 0 + \frac{2}{2}x^2 = 1 + x^2 \le 2, \quad x \in [0,1]
    \]
    即
    \[
        \max_{x \in [0,1]} f(x) \le 2
    \]
\end{proof}
\begin{myRemark}
    利用泰勒公式证明不等式的一般步骤：
    \begin{enumerate}\color {Blue}
        \item 选取适当的展开点$x_0$，一般选取驻点、中点或导数信息多的点．
        \item 选取适当的被展开点$x$，一般选取区间端点．
        \item 确定展开的阶数，能展几阶展几阶．
    \end{enumerate}
\end{myRemark}
\begin{example}
    函数$f(x)$在$[0,1]$上二阶可导，且$f(0)=f(1)=0,\max\limits_{0 \le x \le 1}f(x) = 2$.\\
    证明：$\exists \xi \in (0,1)$，使得$f''(\xi) \le -16$
\end{example}
\begin{proof}
    不妨设$f(x)$在$(0,1)$内的最大值点为$x_0$，则$f(x_0)=2$，且$f'(x_0)=0$．\\
    由泰勒公式知，存在$\xi_1$介于$x_0$与$0$之间，使得
    \[
        f(0)=f(x_0) + f'(x_0)(0 - x_0) + \frac{f''(\xi_1)}{2}(0 - x_0)^2
    \]
    存在$\xi_2$介于$x_0$与$1$之间，使得
    \[
        f(1)=f(x_0) + f'(x_0)(1 - x_0) + \frac{f''(\xi_2)}{2}(1 - x_0)^2
    \]
    因为$f(0)=f(1)=0,f(x_0)=2,f'(x_0)=0$，所以
    \[
        0=2 + \frac{f''(\xi_1)}{2}x_0^2, \quad 0=2 + \frac{f''(\xi_2)}{2}(1 - x_0)^2
    \]
    故
    \[
        f''(\xi_1) = -\frac{4}{x_0^2}, \quad f''(\xi_2) = -\frac{4}{(1 - x_0)^2}
    \]
    因为$x_0 \in (0,1)$，所以
    \[
        \min\left\{f''(\xi_1), f''(\xi_2)\right\} \le -16
    \]
    即存在$\xi \in (0,1)$，使得$f''(\xi) \le -16$．
\end{proof}
\begin{example}
    函数$f(x)$在$[0,1]$上二阶可导，且$|f(x)| \le a, |f''(x)| \le b$，其中$a,b$均为非负常数．
    证明：$\forall x \in (0,1), |f'(x)| \le 2a + \dfrac{b}{2}$
\end{example}
\begin{proof}
    由泰勒公式知，存在$\xi_1$介于$x$与$0$之间，$\xi_2$介于$x$与$1$之间，使得
    \[
        f(0)=f(x) + f'(x)(0 - x) + \frac{f''(\xi_1)}{2}(0 - x)^2
    \]
    \[
        f(1)=f(x) + f'(x)(1 - x) + \frac{f''(\xi_2)}{2}(1 - x)^2
    \]
    于是可得
    \[
        f(1)-f(0)=f'(x)(1 - 0) + \frac{f''(\xi_2)}{2}(1 - x)^2 - \frac{f''(\xi_1)}{2}(0 - x)^2
    \]
    因为$|f(x)| \le a, |f''(x)| \le b$，所以
    \begin{align*}
        |f'(x)|&=\left|f(1)-f(0) - \frac{f''(\xi_2)}{2}(1 - x)^2 + \frac{f''(\xi_1)}{2}x^2\right|\\
        &\le |f(1)| + |f(0)| + \left|\frac{f''(\xi_2)}{2}\right|(1 - x)^2 + \left|\frac{f''(\xi_1)}{2}\right| x^2\\
        &\le a + a + \frac{b}{2}[(1-x)^2 + x^2] = 2a + \frac{b}{2}
    \end{align*}
    其中由于$x \in (0,1)$，故$(1-x)^2 + x^2 \le 1-x+x =1$.
\end{proof}
\ifx\allfiles\undefined
\end{document}
\fi